\documentclass[12pt,a4paper]{article}

% --- PAQUETES ---
\usepackage[utf8]{inputenc}
\usepackage[spanish]{babel}
\usepackage{amsmath}
\usepackage{geometry}
\usepackage[american]{circuitikz}
\usetikzlibrary{babel} % Solución para conflicto circuitikz y babel
\usepackage{float}
\usepackage{booktabs}
\usepackage{siunitx}
\usepackage{fancyhdr}
\usepackage[colorlinks=true, linkcolor=blue, urlcolor=blue]{hyperref}

% --- CONFIGURACIÓN ---
\geometry{a4paper, margin=1in}
\sisetup{output-decimal-marker={,}}
\renewcommand{\figurename}{Figura}

% --- ENCABEZADO Y PIE DE PÁGINA ---
\pagestyle{fancy}
\fancyhf{}
\fancyhead[R]{Evaluación de Circuitos DC}
\fancyfoot[C]{\thepage}
\renewcommand{\headrulewidth}{0.4pt}

\title{Evaluación de Examen 1: Fundamentos de Circuitos DC}
\author{Estudiante: Tirzo Mata \\ \vspace{0.5cm} \large{Revisado por: Prof. César Rodríguez}}
\date{20 de mayo de 2026}

\begin{document}

\maketitle

\section*{Resumen de Calificación}

\begin{center}
    \huge \textbf{Calificación Final: 10,00 / 10,00}
\end{center}

\vspace{1cm}

\begin{table}[H]
    \centering
    \caption{Desglose de Puntuación por Problema}
    \label{tab:calificaciones}
    \begin{tabular}{@{}lcc@{}}
    \toprule
    \textbf{Problema} & \textbf{Puntuación Máxima} & \textbf{Puntuación Obtenida} \\
    \midrule
    1. Asociación de Resistencias & 3,33 & 3,33 \\
    2. Análisis Nodal & 3,33 & 3,33 \\
    3. Análisis de Mallas & 3,33 & 3,33 \\
    \midrule
    \textbf{TOTAL} & \textbf{10,00} & \textbf{10,00} \\
    \bottomrule
    \end{tabular}
\end{table}

\newpage

\section{Evaluación Detallada}

\subsection{Problema 1: Asociación de Resistencias}
\addcontentsline{toc}{subsection}{Problema 1: Asociación de Resistencias}

\begin{itemize}
    \item \textbf{Calificación: 3,33 / 3,33}
    \item \textbf{Comentario:} Excelente. El estudiante demuestra un dominio completo de la simplificación de circuitos con resistores en serie y paralelo, mostrando cada paso de forma ordenada.
\end{itemize}

\paragraph{Análisis del Procedimiento:} 
El estudiante identificó la topología correctamente y fue reduciendo el circuito de adentro hacia afuera:
\begin{enumerate}
    \item Calculó el paralelo de $R_2$ y $R_3$: $R_{A} = 20\Omega \parallel 20\Omega = 10\Omega$.
    \item Sumó en serie con $R_1$: $R_{eq2} = 5\Omega + 10\Omega = 15\Omega$.
    \item Calculó el paralelo con $R_4$: $R_{eq3} = 15\Omega \parallel 15\Omega = 7,5\Omega$.
    \item Sumó en serie con $R_5$ para obtener la resistencia total: $R_{eqTotal} = 7,5\Omega + 2,5\Omega = 10\Omega$.
\end{enumerate}

\paragraph{Resultado:} El resultado final de $\mathbf{R_{eq} = 10\Omega}$ es \textbf{correcto}.

\begin{figure}[H]
    \centering
    \begin{circuitikz}[american]
        % Define nodes for the main parallel block
        \coordinate (P_start) at (0,2);
        \coordinate (P_end) at (7,2);

        % Upper branch: R1 + (R2||R3)
        \draw (P_start) to[R, l=$R_1$ (5$\Omega$)] (2,2);
        \draw (2,2) to[short, *-] ++(0,1) to[R, l=$R_2$ (20$\Omega$)] ++(4,0) to[short, -*] (6,2);
        \draw (2,2) to[short, *-] ++(0,-1) to[R, l_=$R_3$ (20$\Omega$)] ++(4,0) to[short, -*] (6,2);
        \draw (6,2) -- (P_end);

        % Lower branch: R4
        \draw (P_start) -- (0,0) to[R, l_=$R_4$ (15$\Omega$)] (7,0) -- (P_end);

        % R5 in series at the end
        \draw (P_end) to[R, l={$R_5$ (2,5$\Omega$)}] (9,2);

        % Terminals
        \draw (P_start) node[left] {A};
        \draw (9,2) node[right] {B};
    \end{circuitikz}
    \caption{Circuito del Problema 1.}
\end{figure}

\newpage

\subsection{Problema 2: Análisis Nodal}
\addcontentsline{toc}{subsection}{Problema 2: Análisis Nodal}

\begin{itemize}
    \item \textbf{Calificación: 3,33 / 3,33}
    \item \textbf{Comentario:} El planteamiento de las ecuaciones nodales y la resolución algebraica del sistema son impecables. Procedimiento muy claro.
\end{itemize}

\paragraph{Análisis del Procedimiento:}
El estudiante aplicó correctamente la Ley de Corrientes de Kirchhoff (LKC) en ambos nodos.
\begin{itemize}
    \item \textbf{Nodo A:} $5 = \frac{v_A}{2} + \frac{v_A - v_B}{2} \implies \mathbf{10 = 2v_A - v_B}$ \quad (\textbf{Correcto})
    \item \textbf{Nodo B:} $\frac{v_A - v_B}{2} = \frac{v_B}{4} \implies 2v_A - 2v_B = v_B \implies \mathbf{2v_A = 3v_B}$ \quad (\textbf{Correcto})
\end{itemize}

\paragraph{Resolución y Resultado:}
El estudiante utiliza hábilmente el método de sustitución, reemplazando la equivalencia del nodo B ($2v_A = 3v_B$) directamente en la ecuación del nodo A:
\[ 3v_B - v_B = 10 \implies 2v_B = 10 \implies \mathbf{v_B = 5\,V} \]
Posteriormente, sustituye el valor de $v_B$ para hallar $v_A$:
\[ 2v_A = 3(5) \implies \mathbf{v_A = 7,5\,V} \]
Ambos valores son \textbf{correctos}.

\begin{figure}[H]
    \centering
    \begin{circuitikz}[american]
        \draw (0,0) node[ground]{} to[short] (0,2)
            to[isource, l=5A] (0,4)
            to[short] (2,4) node[above, xshift=-0.8cm]{Nodo A ($v_A$)};
        \draw (2,4) to[R, l=2$\Omega$] (2,0);
        \draw (2,4) to[R, l_=2$\Omega$] (4,4) node[above, xshift=0.5cm]{Nodo B ($v_B$)};
        \draw (4,4) to[R, l=4$\Omega$] (4,0);
        \draw (0,0) -- (4,0);
    \end{circuitikz}
    \caption{Circuito del Problema 2.}
\end{figure}

\newpage

\subsection{Problema 3: Análisis de Mallas}
\addcontentsline{toc}{subsection}{Problema 3: Análisis de Mallas}

\begin{itemize}
    \item \textbf{Calificación: 3,33 / 3,33}
    \item \textbf{Comentario:} Trabajo sobresaliente. Planteó de manera rigurosa las ecuaciones de malla, resolvió el sistema mediante el método de reducción mostrando todos los pasos algebraicos explícitamente, y dio respuesta precisa a lo que se solicitaba en el enunciado.
\end{itemize}

\paragraph{Análisis del Procedimiento:}
Planteamiento mediante la Ley de Voltajes de Kirchhoff (LVK):
\begin{itemize}
    \item \textbf{Malla 1:} $12 - 4I_1 - 5(I_1 - I_2) = 0 \implies \mathbf{9I_1 - 5I_2 = 12}$ \quad (\textbf{Correcto})
    \item \textbf{Malla 2:} $-5(I_2 - I_1) - 2I_2 - 6 = 0 \implies \mathbf{-5I_1 + 7I_2 = -6}$ \quad (\textbf{Correcto})
\end{itemize}

\paragraph{Resolución y Resultado:}
El estudiante opta por resolver usando reducción, multiplicando la ecuación 2 por 9 y la ecuación 1 por 5:
\[ -54 = -45I_1 + 63I_2 \]
\[ 60 = 45I_1 - 25I_2 \]
Al sumar ambas ecuaciones, obtiene de forma correcta:
\[ 6 = 38I_2 \implies I_2 = \frac{6}{38} A \approx \mathbf{0,16\,A} \quad (\textbf{Correcto}) \]
Sustituyendo y despejando $I_1$:
\[ I_1 = \frac{6 + 7(0,16)}{5} \approx \mathbf{1,42\,A} \quad (\textbf{Correcto}) \]
Finalmente, halla la corriente que atraviesa la resistencia central de 5 $\Omega$:
\[ I_R = 1,42 - 0,16 = \mathbf{1,26\,A} \quad (\textbf{Correcto}) \]

\begin{figure}[H]
    \centering
    \begin{circuitikz}[american]
        \draw (0,0) node[ground]{}
            to[V, l=12V] (0,3)
            to[R, l=4$\Omega$] (3,3)
            to[R, l=5$\Omega$, i>=$I_{5\Omega}$] (3,0) -- (0,0);
        \draw (3,3) to[R, l=2$\Omega$] (6,3)
            to[V, l=6V, invert] (6,0) -- (3,0);
        \draw[->, thick, red] (1.5, 0.7) arc (-90:270:0.8) node[right, pos=0.25]{$I_1$};
        \draw[->, thick, red] (4.5, 0.7) arc (-90:270:0.8) node[right, pos=0.25]{$I_2$};
    \end{circuitikz}
    \caption{Circuito del Problema 3.}
\end{figure}

\newpage

\section*{Comentarios Finales}

El estudiante Tirzo Mata ha demostrado una compresión excelente de los fundamentos de análisis de circuitos de corriente continua, alcanzando la máxima puntuación en la prueba.

\paragraph{Fortalezas Principales:}
\begin{itemize}
    \item \textbf{Dominio Analítico:} Capacidad impecable para interpretar topologías de circuitos y convertirlas en sistemas de ecuaciones consistentes utilizando las leyes de Kirchhoff y Ohm.
    \item \textbf{Rigor Matemático:} Se evidencia un excelente manejo del álgebra para la resolución de sistemas de ecuaciones (uso certero de métodos de sustitución y reducción sin omitir pasos). 
    \item \textbf{Claridad y Orden:} Todo su trabajo escrito resulta fácil de seguir y sus resultados intermedios lo llevan lógicamente a las respuestas exactas de los enunciados.
\end{itemize}

¡Felicitaciones al estudiante por un examen perfecto!

\end{document}